\section{Implementation}
\label{sec:software}

\subsection{Structure}

The package is organised in three layers. The \texttt{physics} layer
(pure NumPy) provides analytic references and a Numerov integrator used to
cross-check them. The \texttt{solver} layer contains the samplers, networks,
losses and trainer, and---separately from training---the certification
modules \texttt{certify} (1D), \texttt{filtering} (1D DST filter),
\texttt{certify2d} (2D inertia counting and certificates),
\texttt{filtering2d} (2D DST filter and transport) and
\texttt{certify\_subspace} (cluster and invariant-subspace certificates,
Section~\ref{sec:subspace}), \texttt{certified\_count} (directed-rounding
counting and the certified eigenvalue brackets of
Section~\ref{sec:certified}) and \texttt{continuous\_bound} (the variational
bound of Section~\ref{sec:variational}). The \texttt{scripts}
directory contains one driver per experiment plus the aggregation and
reporting utilities; every table in this paper is generated by a script from
saved JSON, not transcribed by hand.

The certification modules do not import PyTorch and do not require a trained
network: they operate on an array of samples and the grid. This is deliberate,
because the certificate must remain meaningful if the training stack changes.

\subsection{Application programming interface}

Table~\ref{tab:api} lists the entry points a user needs; everything else is
built from them. The three certificate calls and the three counting calls are
the whole interface to the numerics.

\begin{table}[t]
\centering
\footnotesize
\setlength{\tabcolsep}{2pt}
\begin{tabular}{@{}ll@{}}
\toprule
Call & Purpose \\
\midrule
\texttt{certify\_eigenpair(psi, x, V, isolate=True)} & 1D certificate \\
\texttt{count\_eigenvalues\_below(x, V, sigma)} & 1D Sturm count \\
\texttt{guided\_filter(psi, x, V)} & 1D filter + re-certification \\
\texttt{certify\_eigenpair\_2d(psi, grid, isolate=True)} & 2D certificate \\
\texttt{count\_eigenvalues\_below\_2d(grid, sigma)} & 2D Sylvester count \\
\texttt{guided\_filter\_2d(psi, grid)} & 2D filter + re-certification \\
\texttt{transported\_angle\_bound\_2d(\dots)} & Lemma~\ref{lem:transport} \\
\texttt{certify\_subspace\_1d(psis, x, V)} & 1D cluster certificate \\
\texttt{certify\_subspace\_2d(psis, grid)} & 2D cluster certificate \\
\texttt{certified\_sturm\_count\_1d(x, V, sigma)} & count, directed rounding \\
\texttt{certified\_separable\_count\_2d(grid, sigma)} & separable 2D count \\
\texttt{certified\_eigenvalue\_brackets\_1d(...)} & certified eigenvalue brackets \\
\texttt{p1\_rayleigh\_bound(x, V, psi)} & upper bound on $\lambda_1(\mathcal H)$ \\
\texttt{q1\_rayleigh\_bound\_2d(x, y, V, psi)} & the same in two dimensions \\
\texttt{principal\_angle(X, exact, h)} & validation only, not a certificate \\
\bottomrule
\end{tabular}
\caption{Main entry points. A certificate is returned as a record with the
energy, the enclosure, the verified gap, the state index, the angle bound and,
when certification is refused, the reason.}
\label{tab:api}
\end{table}

Every certificate is a plain record (a frozen dataclass), so it serialises to
JSON without custom code; the experiment drivers store one record per run. A
refused certificate carries a machine-readable reason---``no isolating
interval exists'', ``gap below resolution floor'', ``Temple condition fails'',
``energy outside the interval''---which is what makes the refusal auditable
rather than a silent \texttt{None}.

A self-contained example, \texttt{examples/certify\_demo.py}, exercises both
paths in about four seconds and without PyTorch: it certifies the analytic trial
function $\mathrm{e}^{-x^2/2}$ of the harmonic oscillator, then reproduces the
refusal-then-subspace sequence of Section~\ref{sec:deg-results} on a
$61\times61$ box, where one member of the degenerate pair is refused as a
state and the pair is certified as a cluster.

\subsection{Drivers and reproducibility}

The experiments of Section~\ref{sec:results} are produced by the commands
below.

\noindent\begin{minipage}{\linewidth}
\footnotesize
\begin{verbatim}
python scripts/run_seed_sweep.py --jobs 6 --threads 4   # 1D, seeds 0-2
python scripts/run_seed_sweep_2d.py --jobs 2            # 2D, seeds 0-2
python scripts/filter_certificates.py                   # 1D certify
python scripts/filter_certificates2d.py                 # 2D certify
python scripts/exp7_degeneracy_2d.py                    # degeneracy
python scripts/exp8_subspace_certificate.py             # clusters
python scripts/exp9_resolution_boundary.py              # floor
python scripts/discretization_study_2d.py --sizes 41 61 81 101 121
python scripts/go_nogo.py ; python scripts/leakage_audit.py
python scripts/multistate_report.py
python scripts/bench_count.py                           # timings
python scripts/dense_reference_2d.py                    # 2D reference
python scripts/appendix_failures.py --variant soft_boundary
python scripts/residual_vs_error.py runs/appendix_failures/unnormalized
python scripts/lapack_reference_widths.py
python scripts/directed_count_report.py
python scripts/continuous_bound.py
python scripts/continuous_bound.py --alignment   # finite well, O(h) sign
python paper/make_tables.py ; python paper/make_figures.py
\end{verbatim}
\end{minipage}

\vspace{0.3em}\noindent Each driver writes a JSON manifest with the wall-clock time of every job, so
the runs are reproducible and the cost figures below are measured rather than
estimated. The 1D runs (10\,000 epochs, grid of 2001 points, network
$64\times3$) take $30$--$45$\,s each on a laptop CPU; the 2D runs
($6000$ epochs, $81\times81$ grid, network $64\times4$) take $99$--$109$\,s
each, manifest times that include start-up and saving.
Certification itself is negligible in both cases. One count
takes $0.4$\,ms at $N=2001$ and $0.034$\,s at $81\times81$ (median of twelve
timed calls, \texttt{scripts/bench\_count.py}); a full certificate---interval
proposal, verified count, and Temple bound---takes $0.074$\,s at $N=2001$,
which is some two hundred counts, and $0.19$\,s at $81\times81$, which is six,
and a complete certificate-guided scan costs $2.6$\,s over the $27$
one-dimensional cutoff orders; in two dimensions the timed benchmark records
$16$\,s for an $81$-candidate scan, and the $30$- and $62$-candidate scans of
the experiments, which examine fewer candidates, scale from it to about $6$
and $12$\,s. The count inherits
the scaling of the sparse factorization it reads: $3$\,ms at $41\times41$,
$0.03$\,s at $81\times81$ and $0.85$\,s at $201\times201$, still small
against a training run on the same grid.

\subsection{Numerical safeguards and tests}

The safeguards described in Section~\ref{sec:counting} are assertions in the
code, not comments: a count that fails the identity-permutation or pivot test
is retried with a perturbed shift and then reported as unavailable. The test
suite (143 tests, run with or without \texttt{pytest}) locks the following
properties, among others:
\begin{itemize}
\item the 1D Sturm count and the 2D Sylvester count agree element by element
  with dense eigen-decomposition on small grids;
\item an exact discrete eigenvector receives a valid certificate whose
  enclosure contains the reference eigenvalue and whose width is at the
  round-off floor;
\item a degenerate eigenvector is refused (Section~\ref{sec:deg-results}),
  while the pair of exact eigenvectors spanning the level is certified as a
  cluster;
\item a rank-deficient set of trial vectors, a single member of a degenerate
  pair, and a cluster whose multiplicity disagrees with the verified count are
  all refused, and a rotated basis of the same subspace returns the same
  certificate;
\item the DST-I used by the filter is its own inverse;
\item the guided filter recovers the exact sine mode of a box potential from a
  perturbed input;
\item the transported bound of Lemma~\ref{lem:transport} dominates the true
  angle;
\item the directed-rounding count agrees with dense eigen-decomposition on
  every shift where it does not refuse, is monotone in the shift, and
  \emph{refuses} at a shift equal to an eigenvalue and for a non-separable
  two-dimensional potential;
\item the variational bound exceeds the true continuous energy for the box
  and the harmonic oscillator, its error contracts at the observed $O(h^2)$
  rate, its quadrature check rejects a kink it was not told about---and, in
  two dimensions, a non-polynomial potential---while accepting the quartic
  anharmonic one, and the two-dimensional Q1 bound equals the sum of the two
  one-dimensional bounds on a non-square grid.
\end{itemize}
The third item deserves emphasis: a test that only asserted ``the certificate
is valid'' would have passed on a version of the code that produced a
spuriously narrow interval for a degenerate state. The test asserts the
refusal.
